\begin{theorem}[exact Bott--Chern and Aeppli diamonds]
\label{thm:exact-bott-chern-aeppli}
\Verified\ The Bott--Chern dimensions of the candidate are
\[
\begin{array}{c|cccc}
q\backslash p&0&1&2&3\\ \hline
3&0&0&1&1\\
2&0&1&2&1\\
1&0&2&1&0\\
0&1&0&0&0,
\end{array}
\tag{11.21}
\]
and its Aeppli dimensions are
\[
\begin{array}{c|cccc}
q\backslash p&0&1&2&3\\ \hline
3&0&0&0&1\\
2&0&1&2&0\\
1&1&2&1&0\\
0&1&1&0&0.
\end{array}
\tag{11.22}
\]
In total degree $k=0,\ldots,6$ this is
\[
 (h^k_{BC})=(1,0,2,2,2,2,1),\qquad
 (h^k_A)=(1,2,2,2,2,0,1).
\tag{11.23}
\]
Consequently, set
\[
 \Delta^k=h^k_{BC}+h^k_A-2b_k
\]
for the degree-$k$ non-$\partial\bar\partial$ defects.  The
Angella--Tomassini inequality
\cite[Theorem~A]{AngellaTomassini2013} gives $\Delta^k\ge0$ for every compact
complex manifold; here the exact values are
\[
 (\Delta^k)=(0,2,4,4,4,2,0).
\tag{11.24}
\]
\end{theorem}

\begin{proof}
McHugh associates to a compact complex manifold the finite-dimensional
quotient
\[
 K^{p,q}=
 \frac{\ker\partial\cap\ker\bar\partial\subset A^{p,q}}
 {\bigl(\operatorname{im}\partial\cap\ker\bar\partial\bigr)
  +\bigl(\operatorname{im}\bar\partial\cap\ker\partial\bigr)},
 \qquad k^{p,q}=\dim_{\mathbb C}K^{p,q}.
\]
Put
\[
 Z^{p,q}=\ker\partial\cap\ker\bar\partial\subset A^{p,q},
 \qquad
 B_K^{p,q}=
 \bigl(\operatorname{im}\partial\cap\ker\bar\partial\bigr)
 +\bigl(\operatorname{im}\bar\partial\cap\ker\partial\bigr).
\]
The natural map $Z^{p,q}\to H^{p+q}_{\mathrm{dR}}(X;\mathbb C)$ need not
descend to an injection from $K^{p,q}$; the dimension estimate has the
opposite quotient mechanism.  If a pure form $\alpha\in Z^{p,q}$ is
$d$-exact, the $(p,q)$ component of a primitive gives
\[
 \alpha=\partial\beta^{p-1,q}+\bar\partial\beta^{p,q-1}.
\]
Because $\bar\partial\alpha=0$ and $\partial\alpha=0$, respectively,
$\partial\beta^{p-1,q}$ is $\bar\partial$-closed and
$\bar\partial\beta^{p,q-1}$ is $\partial$-closed.  Thus
$\ker(Z^{p,q}\to H^{p+q}_{\mathrm{dR}})\subset B_K^{p,q}$, so the de Rham
image of $Z^{p,q}$ surjects onto $K^{p,q}$.  Hence
\[
 k^{p,q}\le b_{p+q}.
\tag{11.25}
\]
Because $X$ has the integral cohomology of $S^6$, (11.25) gives
$k^{p,q}=0$ whenever $0<p+q<6$.

For a compact complex threefold, McHugh's Theorem~6.3,
Proposition~6.4, equation~(7), and the ensuing table give
\begin{align*}
 h^{1,1}_{BC}
 &=2h^{0,1}-b_1+\dim_{\mathbb R}H^{1,1}(X;\mathbb R),\\
 h^{1,2}_{BC}
 &=h^{1,2}+\dim_{\mathbb C}E_2^{3,1}-k^{2,0},\\
 h^{2,2}_{BC}
 &=-h^{1,1}_{BC}+h^{0,1}-h^{0,2}-h^{1,0}+h^{1,1}
   +h^{1,2}+h^{2,0}\\
 &\hspace{2em}+2\dim_{\mathbb C}E_2^{3,1}
   +k^{1,1}-k^{1,2}-2k^{2,0},
\end{align*}
and his threefold table gives the remaining entries
\cite[Theorem~6.3, Proposition~6.4, equation~(7), \S7.1]{McHugh2019}; see also the
independent synthesis in \cite[\S4]{Angella2018}.  Here
$H^{1,1}(X;\mathbb R)$ denotes the subspace of
$H^2_{\mathrm{dR}}(X;\mathbb R)$ represented by real $d$-closed forms of
type $(1,1)$, and it is zero because $H^2_{\mathrm{dR}}(X;\mathbb R)=0$,
while \cref{thm:frolicher-e2-candidate} gives $E_2=E_\infty$; since
$b_4=0$, this in particular gives $E_2^{3,1}=0$.  Substitution of
\[
 h^{0,1}=1,\quad h^{0,2}=h^{1,0}=h^{2,0}=0,\quad
 h^{1,1}=2,\quad h^{1,2}=1
\]
therefore gives
\[
 h^{1,1}_{BC}=2,\qquad h^{1,2}_{BC}=h^{2,1}_{BC}=1,\qquad
 h^{2,2}_{BC}=2.
\]
The same table gives
$h^{2,3}_{BC}=h^{3,2}_{BC}=1$, the two endpoint dimensions equal to one,
and every other entry zero.  This proves (11.21).  The Hodge-star duality
of Schweitzer \cite[\S2.c, pp.~9--10]{Schweitzer2007},
$h_A^{p,q}=h_{BC}^{3-p,3-q}$, proves (11.22).  Summing along diagonals and
using $b_0=b_6=1$ and $b_k=0$ for $1\le k\le5$ proves
(11.23)--(11.24).
\end{proof}
